\documentclass[a4paper,11pt,notitlepage]{report}
% Henrik Kramselund, February 2001
% hlk@security6.net,
% My standard packages
\usepackage{zencurity-exercises}
\begin{document}



\selectlanguage{english}

\newcommand{\subject}
{Secure Development Lifecycle Practices}

%\mytitle{SIEM and Log Analysis}{exercises}
\lhead{\fancyplain{}{\color{titlecolor}\bfseries\LARGE \subject}}


\normal

Revision History:\\
\verb+1.0-create 2026-09-23 HK Created+\\
\verb-1.0-draft 2026-09-29 HK Customer draft+


\section*{Introduction}

This material is prepared for use in \emph{\subject} and was prepared by
Henrik Kramselund, \url{hlk@zencurity.com}.

It contains the plan for the course!

These course and exercises are expected to be performed in a training setting with network connected systems. The exercises use a number of tools which can be copied and reused after training. A lot is described about setting up your workstation in the Git repositories.



\section*{Overall plan}

The overall plan for this course is to outline the idea behind \subject.

\section*{About the course}
\begin{list2}
\item Two days of software security lessons
\item Each day is about 6 lessons each 45min with breaks\\
9:00 - 16:00 with 45min lunch
\item Optional open ended exercises lesson at the end -- 30 min
\item Parts from Software Security course
\end{list2}

Depending on the number of attendees there may be exercises during the day.

Note: this is a very large subject. The main goal is to introduce the most important terms, so the attendees can continue on their journey from an overview of this field.

\section*{Requirements for Secure Development}

A large part of this course is to match requirements for security into the development processes.

Input from ISMS projects

\begin{quote}
ISO 27001:2022, Annex A
\begin{list2}
\item 8.25 Secure development life cycle
\item 8.26 Application security requirements
\item 8.27 Secure system architecture and engineering principles
\item 8.28 Secure coding
\item 8.29 Security testing in development and acceptance
\item 8.30 Outsourced development
\item 8.31 Separation of development, test and production environments
\end{list2}

\end{quote}

Subject plans for the two days, some changes may happen:

Day 0 Basic stuff
\begin{list2}
\item Software security in general – how does things fail
\item History, cases, vulnerabilities, falsehoods and assumptions
\item Security principles: least privilege, fail-safe defaults, etc.
\item OWASP Top 10, CVE/CWE og "Monster mitigations"
\item Security patterns, insecure and secure design
\item Policies, coding standards, compiler warnings, git hooks, defensive programming
\item Code-Auditing Strategies: static analysis and dynamic analysis
\item Security Testing Versus Traditional Software Testing
\item Secure Development Lifecycle (SDL, SDLC, SSDL)
\item Standards ISO27001/CIS Controls, Application Software Security
\end{list2}


Day 1 Advanced subjects and details:
\begin{list2}
\item Authentication: Tokens, MFA, etc.
\item Privacy by Design, Secure by Design
\item Continuous Vulnerability Management
\item Container security
\item Third-party dependencies and examples like Log4Shell
\item Program Building Blocks and Design Patterns
\item Design with Dependencies in Mind, Design for Faults, Design with Promises in Mind
\item Tools and pipelines, Continous Integration - with a few examples
\item Building Secure Infrastructures:
Separation of development, test and production environments, with examples from Kubernetes
\item Enhance and secure runtime environment
\item Logging efficiently -- Maintenance, Monitoring and Analysis of Audit Logs
\end{list2}

Training environments - with Juice Shop as example will be mentioned already on  day 0 and used for demos. This system can easily be setup by developers using either docker og Kubernetes.


\section*{Time schedule}

The above subjects will be put into these lessons:
Plan for the day - day 0
\begin{list2}
\item 09:00 - 09:45 -- 45min TBD
\item 10:00 - 10:45 -- 45min TBD
\item 11:00 - 11:45 -- 45min TBD
\item 11:45 - 12:30 -- Lunch
\item 12:30 - 13:15 -- 45min TBD
\item 13:30 - 14:15 -- 45min TBD
\item 14:30 - 15:15 -- 45min TBD
\item 15:30 - 16:00 -- 30min Optional: Exercises, questions
\end{list2}

Plan for the day - day 1
\begin{list2}
\item 09:00 - 09:45 -- 45min TBD
\item 10:00 - 10:45 -- 45min TBD
\item 11:00 - 11:45 -- 45min TBD
\item 11:45 - 12:30 -- Lunch
\item 12:30 - 13:15 -- 45min TBD
\item 13:30 - 14:15 -- 45min TBD
\item 14:30 - 15:15 -- 45min TBD
\item 15:30 - 16:00 -- 30min Optional: Exercises, questions
\end{list2}

\end{document}
